\documentclass[12pt]{extarticle}
\usepackage[a4paper,left=2.8cm,right=2.8cm,top=2.2cm,bottom=2cm]{geometry}
\usepackage{amsmath,xcolor,enumitem,amssymb,qrcode}
\pagestyle{empty}
\setlength{\parindent}{0pt}
\setlist{leftmargin=1.4em,itemsep=3pt,topsep=4pt}
\newcommand{\abschnitt}[1]{\par\vspace{22pt}{\large\bfseries #1}\par\vspace{5pt}}
\newcommand{\jj}{\,$\bullet$}
\begin{document}
{\LARGE\bfseries MSA · P10 Mathematik}\hfill{\small\color{gray}Stand Oktober 2026}
\abschnitt{Termin}
Dienstag, 4.\,Mai 2027, 135 Minuten
\abschnitt{Zwei Prüfungen}
\begin{tabular}{@{\hspace{1.4em}}l@{\hspace{1em}}c@{\hspace{1em}}l@{}}
Grundkurs & $\longrightarrow$ & \textbf{EBR-Prüfung}\\[3pt]
Erweiterungskurs & $\longrightarrow$ & \textbf{FOR-Prüfung}\\
\end{tabular}
\par\vspace{8pt}
\textbf{Basisaufgaben} (Aufgabe 1) sind in beiden gleich.
\abschnitt{Mitbringen}
Taschenrechner, Formelsammlung, Geodreieck, Zirkel
\abschnitt{Was am meisten Punkte bringt}
{\small Anteil an allen Punkten der Prüfungen 2020–2026}\par\vspace{6pt}
\begin{minipage}{10cm}
\begin{enumerate}[label=\arabic*.,leftmargin=1.8em,itemsep=4pt]
\item Lineare Funktionen \hfill 11\,\%\jj
\item Quadratische Funktionen \hfill 10\,\%\jj
\item Mittelwert, Median, Spannweite \hfill 8\,\%\jj
\item Wahrscheinlichkeit \hfill 7\,\%\phantom{\jj}
\item Flächen und Umfang \hfill 6\,\%\phantom{\jj}
\item Volumen und Oberfläche \hfill 6\,\%\phantom{\jj}
\item Satz des Pythagoras \hfill 6\,\%\phantom{\jj}
\item Trigonometrie \hfill 6\,\%\jj
\item Wachstum \hfill 5\,\%\phantom{\jj}
\item Prozentrechnung \hfill 4\,\%\jj
\end{enumerate}
\end{minipage}
\par\vspace{6pt}
{\small $\bullet$ kommt fast jedes Jahr dran}
\abschnitt{Alte Prüfungen zum Üben}
\begin{minipage}[c]{0.68\linewidth}
Alle Prüfungen seit 2014, frei zum Herunterladen:\par\vspace{4pt}
{\small\raggedright Bildungsserver Berlin-Brandenburg $\to$ Prüfungen $\to$ Prüfungen~10 $\to$ Prüfungsaufgaben Mathematik}
\end{minipage}\hfill
\begin{minipage}[c]{0.26\linewidth}\raggedleft
\qrcode[height=2.6cm]{https://bildungsserver.berlin-brandenburg.de/unterricht/pruefungen/pruefungen-10/pruefungsaufgaben-mathematik}
\end{minipage}
\end{document}
